\chapter{Introduzione}

\section{Scopo della Tesi}
Lo scopo principale di questo elaborato è indagare empiricamente e analiticamente la robustezza dei moderni Large Language Models applicati al task di Vulnerability Detection nel codice sorgente C. Piuttosto che limitarsi a misurare l'accuratezza superficiale dei modelli, la ricerca si propone di stressare le architetture neurali attraverso attacchi avversariali mirati, valutando come perturbazioni testuali e strutturali possano indurre i modelli in errore, così da essere sfruttati da attaccanti per eseguire codice malevolo. 

L'elaborato mira a superare il tradizionale approccio puramente osservazionale, integrando tecniche avanzate di Mechanistic Interpretability e Representation Engineering. Lo scopo ultimo è tracciare la traiettoria computazionale dei LLM per comprendere le dinamiche latenti che guidano la classificazione del codice, dimostrando geometricamente come il modello codifichi nello spazio latente le perturbazioni introdotte dagli attacchi avversari.

\section{Contesto e Motivazione}
Nella società contemporanea, data la sempre più crescente presenza dei sistemi informatici, il software rappresenta una componente fondamentale e potenzialmente critica per settori ed enti importanti come banche, servizi sul territorio come telecomunicazioni, sanità e gestione dei dati personali.

Ad un incremento delle necessità e dell'impiego di software parallelamente è cresciuta la complessità del codice sorgente e l'area di attacco esposta a potenziali minacce.
Le vulnerabilità del codice, soprattutto per linguaggi a basso livello come C o C++, rappresentano una delle principali sfide per la sicurezza informatica causando possibili danni economici e reputazionali.
Storicamente per rilevare possibili minacce sono stati usati strumenti di analisi statica e dinamica, o addirittura semplice ricerche testuali, per cotnrastare possibili falle nel codice.
Tuttavia questo approccio hanno dimostrato diversi limiti strutturali: da un lato, l'analisi statica produce diversi falsi positivi dall'altro le continue revisioni manuali da parte degli sviluppatori risulta insostenibile di fronte alla mole di codice generata nei moderni cicli di sviluppo software.

In questo scenario, l'avvento dell'Intelligenza Artificiale Generativa e, in particolare, dei Large Language Models (LLM) ha inaugurato un cambio di paradigma. Modelli addestrati su vaste linee di codice sorgente hanno dimostrato capacità sorprendenti non solo nella generazione di codice, ma anche nell'analisi e nell'identificazione di pattern complessi. L'adozione degli LLM come "assistenti alla sicurezza" promette di superare i limiti semantici degli analizzatori tradizionali.

Tuttavia, l'integrazione di questi modelli in casi di sicurezza critiche solleva interrogativi urgenti sulla loro affidabilità e robustezza. Mentre le prestazioni degli LLM vengono tipicamente valutate in scenari statici in condizioni ideali, ma il dominio della cybersecurity è per natura applicata in un ambiente ostile. Un attaccante può tentare di ingannare il modello manipolando il codice vulnerabile affinché appaia sicuro, aggirando così le difese automatizzate. Comprendere quanto questi modelli siano resilienti agli attacchi e, soprattutto, perché falliscano nel ragionamento logico, rappresenta oggi una priorità assoluta per la ricerca nell'ambito della sicurezza dell'Intelligenza Artificiale.

\section{Sintesi dei contenuti}

\begin{itemize}
    \item \textbf{Capitolo 2 - Backgroud e Stato dell'Arte:} In questo capitolo vengono introdotti i concetti fondamentali necessari per comprendere il contesto in cui si inserisce il presente lavoro di tesi. La trattazione si articola su tre assi principali: l'evoluzione e l'architettura dei Large Language Models (LLM), con particolare attenzione al meccanismo dei Transformer e allo spazio latente; il ruolo di tali modelli nell'ambito della sicurezza del software per la vulnerability detection; e infine, la tassonomia degli attacchi avversariali, supportata dai principi della Mechanistic Interpretability.
    \item \textbf{Capitolo 3 - Metodologia adottata:} 
    \item \textbf{Capitolo 4 - Analisi delle Performance:}
    \item \textbf{Capitolo 5 - Operazioni di Mechanistic Interpretability:}
    \item \textbf{Capitolo 6 - Representation Engineering:}
    \item \textbf{Capitolo 7 - Conclusioni}
\end{itemize}









